\documentclass[10pt,a4paper]{amsart}
\usepackage[margin=19mm]{geometry}
\usepackage{amsmath,amssymb,mathtools}
\usepackage[scr=rsfs,cal=euler]{mathalfa}
\usepackage{xfrac}
\usepackage[unicode,hidelinks,bookmarks=false]{hyperref}
\newcommand{\ie}{i.e.\ }
\newcommand{\cf}{cf.\ }
\newcommand{\resp}{respectively\ }
\newcommand{\Subsection}{\subsection}
\providecommand{\nobreakdash}{\nobreak\hbox{-}}
\setlength{\emergencystretch}{2em}
\multlinegap0pt
\usepackage{bm}
\usepackage{bbm}
\usepackage{dsfont}
\usepackage{xspace}
\usepackage{cancel}
\usepackage{mathtools}
\usepackage{amsthm}
\makeatletter
\@ifundefined{lemma}{\newtheorem{lemma}{Lemma}}{}
\renewenvironment{proof}[1]{\par\medskip\noindent\textbf{#1}\enspace\ignorespaces}{\par\medskip}
\makeatother
\providecommand{\blambda}{\boldsymbol{\lambda}}
\providecommand{\bzero}{\boldsymbol{0}}
\title[Dual-Based Weight Selection for Approximate Linear Programmi]{Dual-Based Weight Selection for Approximate Linear Programming}
\author{Su Li, Andre A. Cire, Adam Diamant, Vahid Sarhangian}
\date{Source: arXiv:2608.24629v1 (CC BY 4.0)}
\begin{document}
\maketitle

\noindent\emph{Source and licence.} Dual-Based Weight Selection for Approximate Linear Programming. Version arXiv:2608.24629v1 (CC BY 4.0), distributed under the Creative Commons Attribution 4.0 International licence, as recorded in the arXiv licence metadata for this exact version and shown on the abstract page https://arxiv.org/abs/2608.24629v1. Full rights notice: licenses/arxiv-2608.24629-CC-BY-4.0.txt.\par\smallskip

\noindent\textit{Editorial scope.} Counted unit: the Lemma whose source label is \texttt{lemma:L\_Lipschitz\_policy} in the article (source0.tex lines 609--628, source label \texttt{lemma:L\_Lipschitz\_policy}), delivered together with the article's own complete original proof of it (source0.tex lines 629--706, one \texttt{proof} environment of 78 lines). No mathematics is written, repaired or re-derived: statement and proof are the article's own text line for line. Required context retained uncounted: eq:alp\_dual\_expanded (lines 374--397); eq:target\_guided\_softmax\_policy (lines 448--462). Every definition, display and label of those lines is retained unchanged. Omitted from this package and not redistributed: 1--373, 398--447, 463--608, 707--2959. Nothing inside a retained excerpt is omitted or altered: every retained body is delivered exactly as the article's lines are written, so no omission receipt of any kind accompanies this package. Citations: the retained excerpts carry no citation command at all, so no bibliography entry of the article is transcribed and no reference is redistributed. Reference adaptation: none was needed and none was made; every reference inside the delivered unit resolves inside it, so no unresolved reference or citation is produced. Printed slips kept literal: no printed word, symbol, hypothesis or number of the counted statement or of its proof was altered. Layout and class: the article's own class is \texttt{informs4} with packages including fix-cm, tgtermes, newtxtext, newtxmath, bm, endnotes, microtype, xcolor, graphicx, amsmath, algorithm2e, acronym, comment, booktabs; the wrapper is the lane's pinned \texttt{amsart} editorial wrapper at 10pt on A4, which restates the theorem-like environments the retained text needs and adds the article's own macro definitions that the retained text needs (\texttt{\textbackslash blambda}, \texttt{\textbackslash bzero}), with extra wrapper packages (bm, bbm, dsfont, xspace, cancel, mathtools). The delivered numbering is the unit's own. Assets: no retained span contains a figure environment, a graphics-inclusion command, a PDF-page inclusion, a table or a TikZ picture; no image file is copied into this package, the package declares no asset dependency and no image file is redistributed with it. Licence: version arXiv:2608.24629v1 of this article is distributed under the Creative Commons Attribution 4.0 International licence, as recorded in the arXiv licence metadata for this exact version and shown on the abstract page https://arxiv.org/abs/2608.24629v1. A full rights notice is delivered at licenses/arxiv-2608.24629-CC-BY-4.0.txt. Characters: every retained excerpt is pure ASCII, so the delivered engine, XeLaTeX with its default Latin Modern text font, reports no missing character.
\begin{equation}
\label{eq:alp_dual_expanded}
\tag{dALP}
\begin{aligned}
  \min_{\blambda\geq\bzero}\quad
  &
  \sum_{x\in\mathcal S}
  \sum_{a\in\mathcal A(x)}
  g(x,a)\lambda(x,a) \\
  \text{s.t.}\quad
  &
  \sum_{x\in\mathcal S}
  \sum_{a\in\mathcal A(x)}
  \lambda(x,a)
  \left[
  \phi_k(x)
  -
  \alpha\sum_{y\in\mathcal S}P(y\mid x,a)\phi_k(y)
  \right]
  =
  \sum_{x\in\mathcal S}c(x)\phi_k(x),
  \;\; \forall k=1,\ldots,K.
\end{aligned}
\end{equation}

\begin{equation}
\label{eq:target_guided_softmax_policy}
\begin{aligned}
\pi_{\tau,\epsilon,q}(a\mid x;\overline{\blambda})
&=
\frac{\sum_{b\in\mathcal A(x)}\overline{\lambda}(x,b)}
     {\sum_{b\in\mathcal A(x)}\overline{\lambda}(x,b)+\epsilon}
\frac{\exp\!\left(\overline{\lambda}(x,a)/\tau\right)}
     {\sum_{a'\in\mathcal A(x)}\exp\!\left(\overline{\lambda}(x,a')/\tau\right)}
+
\frac{\epsilon}
     {\sum_{b\in\mathcal A(x)}\overline{\lambda}(x,b)+\epsilon}
q(a\mid x).
\end{aligned}
\end{equation}

\begin{lemma}[Policy Stability.]
\label{lemma:L_Lipschitz_policy}
There exist parameters $\epsilon, \tau > 0$ such that, for any two dual feasible solutions $\blambda, \blambda'$ to \eqref{eq:alp_dual_expanded} and a prespecified target policy $q$,
\[
\lVert \pi_{\tau,\epsilon,q}(\cdot \mid x;\blambda) - \pi_{\tau,\epsilon,q}(\cdot\mid x; \blambda') \rVert_1\le \, L_p \lVert \blambda_x - \blambda'_x \rVert_1, 
\;\; \forall x \in\mathcal S, 
\]
for some constant $L_p < (1-\alpha)/(L_d\alpha)$. Specifically, it suffices to pick any $\epsilon, \tau > 0$ such that  
\begin{equation}
\label{eq:epsilon_tau_conditions}
  \epsilon
  >
  \frac{2\alpha L_d}{1-\alpha},
  \qquad
  \tau
  >
  \frac{\alpha L_d\epsilon}
  {2\left[(1-\alpha)\epsilon-2\alpha L_d\right]}.
\end{equation}
\end{lemma}

\begin{proof}{Proof of Lemma \ref{lemma:L_Lipschitz_policy}.}
Let $x\in\mathcal S$ and $q_x=\bigl(q(a\mid x)\bigr)_{a\in\mathcal A(x)}$. For any $\blambda_x \ge 0$, define
\begin{equation*}
\sigma_\tau(a;\blambda_x)
=
\frac{\exp\!\left(\lambda(x,a)/\tau\right)}
{\sum_{b\in\mathcal A(x)}\exp\!\left(\lambda(x,b)/\tau\right)},
\quad
s(\blambda_x)=\sum_{a\in\mathcal A(x)}\lambda(x,a),
\quad
\omega_\epsilon(\blambda_x)=\frac{s(\blambda_x)}{s(\blambda_x)+\epsilon}.
\end{equation*}
Then, the smoothed policy \eqref{eq:target_guided_softmax_policy} becomes
$\pi_{\tau,\epsilon,q}(\cdot\mid x;\blambda)
=\omega_\epsilon(\blambda_x)\sigma_\tau(\blambda_x)
+[1-\omega_\epsilon(\blambda_x)]q_x$.

For any two actions $a,b\in\mathcal A(x)$, the derivative of the softmax term yields
\begin{equation*}
\frac{\partial\sigma_\tau(a;\blambda_x)}
{\partial\lambda(x,b)}
=
\frac{1}{\tau}\sigma_\tau(a;\blambda_x)
\bigl[\mathbf 1\{a=b\}-\sigma_\tau(b;\blambda_x)\bigr].
\end{equation*}
For each column $b$ of the Jacobian,
\begin{equation*}
\sum_{a\in\mathcal A(x)}
\left|
\frac{\partial\sigma_\tau(a;\blambda_x)}
{\partial\lambda(x,b)}
\right|
=
\frac{2}{\tau}\sigma_\tau(b;\blambda_x)
\bigl[1-\sigma_\tau(b;\blambda_x)\bigr]
\leq
\frac{1}{2\tau}.
\end{equation*}
The equality uses
$\sum_{a\neq b}\sigma_\tau(a;\blambda_x)
=1-\sigma_\tau(b;\blambda_x)$, and the inequality follows from
$p(1-p)\leq1/4$ for any scalar $p\in[0,1]$. Hence, the induced matrix $1$-norm of
the softmax Jacobian is at most $1/(2\tau)$. Applying the integral form
of the mean-value theorem along the line segment between $\blambda_x$
and $\blambda'_x$ therefore gives
$\|\sigma_\tau(\blambda_x)-\sigma_\tau(\blambda'_x)\|_1
\leq\|\blambda_x-\blambda'_x\|_1/(2\tau)$.

Similarly, the derivative of the ratio $s/(s+\epsilon)$ for any scalar $s\geq0$ is bounded by
$1/\epsilon$. Since
$|s(\blambda_x)-s(\blambda'_x)|
\leq\|\blambda_x-\blambda'_x\|_1$, it follows that
$|\omega_\epsilon(\blambda_x)-\omega_\epsilon(\blambda'_x)|
\leq\|\blambda_x-\blambda'_x\|_1/\epsilon$.

Let $\omega=\omega_\epsilon(\blambda_x)$,
$\omega'=\omega_\epsilon(\blambda'_x)$,
$\sigma=\sigma_\tau(\blambda_x)$, and
$\sigma'=\sigma_\tau(\blambda'_x)$. Since
$\pi(\blambda_x)-\pi(\blambda'_x)
=\omega(\sigma-\sigma')+(\omega-\omega')(\sigma'-q_x)$,
$0\leq\omega\leq1$, and $\|\sigma'-q_x\|_1\leq2$, the preceding
bounds yield
\begin{equation*}
\left\|
\pi_{\tau,\epsilon,q}(\cdot\mid x;\blambda)
-\pi_{\tau,\epsilon,q}(\cdot\mid x;\blambda')
\right\|_1
\leq
\left(\frac{1}{2\tau}+\frac{2}{\epsilon}\right)
\left\|\blambda_x-\blambda'_x\right\|_1.
\end{equation*}
Thus, we may take $L_p=1/(2\tau)+2/\epsilon$. Under
\eqref{eq:epsilon_tau_conditions},
$1/(2\tau)< (1-\alpha)/(\alpha L_d)-2/\epsilon$, and hence
$L_p<(1-\alpha)/(\alpha L_d)$, as required.
$\hfill \blacksquare$
\end{proof}

\end{document}
